% ==========================================================
% 圆排列中的插空
% 难度：★★★ 难题
% ==========================================================

\newcommand{\qProbPermCircleInsertStem}{%
五个女孩与六个男孩围成一圈，要求任意两个女孩中间至少站一个男孩。求不同排法数。所有人均视为互不相同，圆排列中旋转相同的排列视为同一种。
}

\newcommand{\qProbPermCircleInsertAnswer}{%
$86400$
}

\newcommand{\qProbPermCircleInsertSolution}{%
\textbf{解法一（女孩先排）：}\\
先把五个女孩围成一圈，\\
圆排列数为 $(5-1)!=4!$。\\
五个女孩之间形成五个圆形空隙。\\
为使两个女孩间至少有一个男孩，\\
这五个空隙都必须有男孩。\\
六名男孩分入五个非空空隙，\\
人数结构只能是 $(2,1,1,1,1)$。\\
选择哪两名男孩组成双人组：$C_6^2$。\\
双人组在同一空隙中有 $2!$ 种先后顺序。\\
此时形成五个不同的男孩单位。\\
把五个单位安排到女孩形成的五个固定空隙中，\\
由于空隙的位置已相对固定，\\
这是普通的直线排列：$5!$ 种。\\
总排法数 $N = 4!\cdot C_6^2\cdot 2!\cdot 5! = 86400$ 种。

\textbf{解法二（男孩先排更简洁）：}\\
先把六名男孩围成一圈，\\
圆排列数为 $(6-1)!=5!$。\\
六名男孩形成六个圆形空隙。\\
从六个空隙中选择五个放入五名女孩，\\
方法数为 $A_6^5$。\\
总排法数 $N = 5!\cdot A_6^5 = 120\cdot 720 = 86400$。
}

\newcommand{\qProbPermCircleInsertTeachingNotes}{%
\textbf{【避坑与边界说明】}
\begin{itemize}
    \item 线性排列中，$n$ 个对象产生 $n+1$ 个空隙；\\
    圆排列中，$n$ 个对象只产生 $n$ 个空隙。
    \item 圆排列固定一类对象后，空隙已有标号。\\
    后续插空不能再使用圆排列公式，而是直线排列。
    \item 男孩先排方法更接近普通插空法，\\
    且不需要讨论人数结构，计算更为简捷。
\end{itemize}
}
